\subsection{Kummer theory}

In this No.~we denote by n an integer \textgreater0, and we shall suppose that \(\mu_{n}(\mathbf{K})\) has n elements; by V, p.~78 this means also that n is prime to the characteristic exponent of K and that all the n-th roots of unity in an algebraic closure \(\Omega\) of K lie in K.

We shall say that an extension L of K is abelian of exponent dividing n if it is abelian (V, p.~77, Def. 1) and its Galois group \(\operatorname{Gal}(\mathrm{L}/\mathrm{K})\) is annihilated by n (V, p.~86).

Let A be a subset of \(K^{*}\) ; we denote by \(\mathbf{K}(\mathbf{A}^{1/n})\) the subextension of \(\Omega\) generated by all \(\theta \in \Omega\) such that \(\theta^{n} \in A\) .

Lemma 5. --- \(\mathbf{K}(\mathbf{A}^{1 / n})\) is an abelian extension of \(\mathbf{K}\) of exponent dividing \(n\) .

Since the polynomials \(X^n - a\) , \(a \in A\) are separable over \(K\) , \(L = K(A^{1/n})\) is a separable, hence a Galois extension of \(K\) . Let \(\sigma \in \operatorname{Gal}(L/K)\) and let \(\theta \in \Omega\) be such that \(\theta^n \in A\) . We have \(\sigma(\theta)^n = \theta^n\) ; hence there exists \(\zeta \in \mu_n(\Omega) = \mu_n(K)\) such that \(\sigma(\theta) = \zeta\theta\) ; this implies that \(\sigma^n(\theta) = \zeta^n\theta = \theta\) , whence \(\sigma^n = 1\) . If \(\sigma'\) is another element of \(\operatorname{Gal}(L/K)\) there exists \(\zeta' \in \mu_n(K)\) such that \(\sigma'(\theta) = \zeta'\theta\) , whence \(\sigma' \sigma(\theta) = \zeta\zeta'\theta = \sigma\sigma'(\theta)\) and so \(\sigma' \sigma = \sigma\sigma'\) .

Lemma 6. --- Let L be a Galois extension of K. There exists a unique mapping \((\sigma, a) \mapsto \langle \sigma, a \rangle\) of \(\operatorname{Gal}(\mathrm{L} / \mathrm{K}) \times ((\mathrm{L}^n \cap \mathrm{K}^*) / \mathrm{K}^{*n})\) into \(\mu_n(\mathbf{K})\) such that for every \(\sigma \in \operatorname{Gal}(\mathrm{L} / \mathrm{K})\) and every element \(\theta \in \mathrm{L}^*\) such that \(\theta^n \in \mathbf{K}\) we have, on denoting by \(\bar{\theta}^n\) the residue class of \(\theta^n \bmod \mathrm{K}^{*n}\) :

\[
\left\langle \sigma , \bar {\theta} ^ {n} \right\rangle = \sigma (\theta) / \theta .\tag{13}
\]

This mapping is bimultiplicative.

In effect the right-hand side of (13) is an n-th root of unity which only depends on the residue class mod \(K^{*n}\) of \(\theta^{n}\) ; this proves the first assertion. The second assertion is verified without difficulty.

For every Galois extension L of K let us write

\[
\begin{array}{l} k _ {\mathrm{L}} \colon (\mathrm{L} ^ {n} \cap \mathrm{K} ^ {*}) / \mathrm{K} ^ {* n} \to \operatorname{Hom} \left(\operatorname{Gal} (\mathrm{L} / \mathrm{K}),   \mu_ {n} (\mathrm{K})\right), \\ k _ {\mathrm{L}} ^ {\prime} \colon \operatorname{Gal} (\mathrm{L} / \mathrm{K}) \to \operatorname{Hom} \left((\mathrm{L} ^ {n} \cap \mathrm{K} ^ {*}) / \mathrm{K} ^ {* n},   \mu_ {n} (\mathrm{K})\right), \end{array}
\]

for the homomorphisms derived from the above bimultiplicative mapping (V, p.~87).

PROPOSITION 9. --- For every Galois extension of finite degree L of K, the homomorphism \(k_{L}\) is bijective.

Let \(\theta \in L^{*}\) be such that \(\theta^n \in K\) and the residue class of \(\theta^n \bmod K^{*n}\) belongs to the kernel of \(k_{L}\) . For every \(\sigma \in \operatorname{Gal}(L/K)\) we have by definition \(\sigma(\theta) = \theta\) ; hence \(\theta \in K^{*}\) and \(\theta^n \in K^{*n}\) . This proves the injectivity of \(k_{L}\) . Now let \(f: \operatorname{Gal}(L/K) \to \mu_n(K)\) be a homomorphism; for all \(\sigma, \tau \in \operatorname{Gal}(L/K)\) we have

\[
f (\sigma \tau) = f (\sigma) f (\tau) = f (\sigma). \sigma f (\tau), \quad f (\sigma) ^ {n} = 1.
\]

By V, p.~65, Cor. 1 there exists \(\theta \in L^{*}\) such that \(f(\sigma) = \sigma(\theta)/\theta\) for all \(\sigma \in \operatorname{Gal}(L/K)\) ; since \(f(\sigma)^{n} = 1\) , we have \(\sigma(\theta^{n}) = \theta^{n}\) for all \(\sigma \in \operatorname{Gal}(L/K)\) , hence \(\theta^{n} \in K^{*}\) ; if \(a\) is the residue class of \(\theta^{n} \bmod K^{*n}\) , we have by definition \(f(\sigma) = \langle \sigma, a \rangle\) for \(\sigma \in \operatorname{Gal}(L/K)\) , hence \(f = k_{L}(a)\) .

COROLLARY. --- If L is a Galois extension of K, the homomorphism \(k_{L}\) is injective and its image is the group \(\operatorname{Hom}_{c}(\operatorname{Gal}(\operatorname{L}/\operatorname{K}), \mu_{n}(\operatorname{K}))\) of continuous homomorphisms of the topological group \(\operatorname{Gal}(\operatorname{L}/\operatorname{K})\) into the discrete group \(\mu_{n}(\operatorname{K})\) .

This follows at once from what has been said, using the fact that L is an increasing directed union of Galois extensions \(L_{i}\) of finite degree and that the homomorphism of \(\operatorname{Gal}(L/K)\) into \(\mu_{n}(K)\) is continuous if and only if it can be factored through one of the quotients \(\operatorname{Gal}(L_{i}/K)\) of \(\operatorname{Gal}(L/K)\) .

THEOREM 4. --- a) The mapping \(\mathrm{H} \mapsto \mathrm{K}\left(\mathrm{H}^{1/n}\right)\) is an inclusion preserving bijection of the set of subgroups of \(\mathbf{K}^*\) containing \(\mathbf{K}^{*n}\) onto the set of abelian subextensions of exponent dividing \(n\) of \(\Omega\) . The inverse mapping is \(\mathrm{L} \mapsto \mathrm{L}^n \cap \mathbf{K}^*\) .

\begin{enumerate}
\def\labelenumi{\alph{enumi})}
\setcounter{enumi}{1}
\tightlist
\item
  For every subgroup H of \(\mathbf{K}^*\) containing \(\mathbf{K}^{*n}\) the homomorphism
\end{enumerate}

\[
k ^ {\prime}: \operatorname{Gal} \left(\mathrm{K} \left(\mathrm{H} ^ {1 / n}\right) / \mathrm{K}\right)\rightarrow \operatorname{Hom} \left(\mathrm{H} / \mathrm{K} ^ {* n}, \mu_ {n} (\mathrm{K})\right)
\]

is bijective, and it is a homeomorphism when the group \(\operatorname{Hom}(\mathrm{H} / \mathrm{K}^{*n},\mu_n(\mathrm{K}))\) is equipped with the topology of simple convergence.

\begin{enumerate}
\def\labelenumi{\alph{enumi})}
\setcounter{enumi}{2}
\tightlist
\item
  Let H be a subgroup of \(K^{*}\) containing \(K^{*n}\) . For every \(a \in H/K^{*n}\) let \(\theta_{a}\) be an element of \(\Omega\) such that \(\theta_{a}^{n}\) is a representative of a in H. Then the \(\theta_{a}\) , \(a \in H/K^{*n}\) , form a basis of the vector K-space \(\mathrm{K}(H^{1/n})\) . In particular,
\end{enumerate}

\[
[ \mathrm{K} (\mathrm{H} ^ {1 / n}): \mathrm{K} ] = (\mathrm{H}: \mathrm{K} ^ {* n}).
\]

\begin{enumerate}
\def\labelenumi{\Alph{enumi})}
\tightlist
\item
  For every abelian extension L of K of exponent dividing n let us put \(H_{L} = L^{n} \cap K^{*}\) . If \([L : K]\) is finite, the homomorphism \(k_{L}'\) of \(\operatorname{Gal}(L/K)\) into \(\operatorname{Hom}(H_{L}, \mu_{n}(K))\) is bijective by Prop. 9 and V, p.~88, Prop. 8.
\end{enumerate}

Since every abelian extension of K of exponent dividing n is an increasing directed union of abelian subextensions of finite degree of exponent dividing n, we deduce by passage to inverse limits that \(k_{L}^{\prime}\) is a homeomorphism of topological groups for every abelian extension L of K dividing n.

\begin{enumerate}
\def\labelenumi{\Alph{enumi})}
\setcounter{enumi}{1}
\item
  Let L be an abelian extension of finite degree, of exponent dividing n, of K, and let \(L' = \mathbf{K}(\mathbf{H}_{L}'^{1/n})\) ; this is a subextension of L; moreover \(H_{L'}\) contains \(H_{L}\) and so is equal to it. Since the homomorphisms \(k_{L}'\) and \(k_{L'}\) are bijective by A) and \(H_{L} = H_{L'}\) , the groups \(\operatorname{Gal}(L/K)\) and \(\operatorname{Gal}(L'/K)\) have the same order and so are equal. This shows that \(L' = L\) , and hence that \(L = \mathbf{K}(\mathbf{H}_{L}'^{1/n})\) . If L is an abelian extension of exponent dividing n of K, we have \(\mathbf{K}(\mathbf{H}_{L}'^{1/n}) = \mathbf{L}\) , because \(\mathbf{K}(\mathbf{H}_{L}'^{1/n})\) is a subextension of L which contains every subextension of finite degree of L.
\item
  Let H be a subgroup of \(K^{*}\) containing \(K^{*n}\) ; put \(L = K(H^{1/n})\) , then this is an abelian extension of K of exponent dividing n (V, p.~88, Lemma 5). We have \(H \subset H_{L}\) , whence we obtain an exact sequence of commutative groups annihilated by n
\end{enumerate}

\[
\{1 \} \rightarrow \mathrm{H} / \mathrm{K} ^ {* n} \rightarrow \mathrm{H} _ {\mathrm{L}} / \mathrm{K} ^ {* n} \rightarrow \mathrm{H} _ {\mathrm{L}} / \mathrm{H} \rightarrow \{1 \}.
\]

From this we obtain an exact sequence

\[
\left\{1 \right\}\rightarrow \mathrm{Hom} \left(\mathrm{H} _ {\mathrm{L}} / \mathrm{H}, \mu_ {n} (\mathrm{K})\right)\rightarrow \mathrm{Hom} \left(\mathrm{H} _ {\mathrm{L}} / \mathrm{K} ^ {* n}, \mu_ {n} (\mathrm{K})\right) \stackrel {u} {\rightarrow} \mathrm{Hom} \left(\mathrm{H} / \mathrm{K} ^ {* n}, \mu_ {n} (\mathrm{K})\right)
\]

where \(u\) is the restriction homomorphism.

If we identify \(\operatorname{Hom}(\mathrm{H}_{\mathrm{L}}/\mathrm{K}^{*n},\mu_{n}(\mathrm{K}))\) with \(\operatorname{Gal}(\mathrm{L}/\mathrm{K})\) by means of the isomorphism \(k_{L}^{\prime}\) , the kernel of u is identified with the set of \(\sigma\in\operatorname{Gal}(\mathrm{L}/\mathrm{K})\) such that \(\sigma(\theta)=\theta\) for all \(\theta\in H^{1/n}\) . It follows that u is injective, hence \(\operatorname{Hom}(\mathrm{H}_{\mathrm{L}}/\mathrm{H},\mu_{n}(\mathrm{K}))\) is reduced to \(\{1\}\) ; by Cor. 1 of V, p.~87 we thus have \(H=H_{L}\) . This completes the proof of a) and b).

\begin{enumerate}
\def\labelenumi{\Alph{enumi})}
\setcounter{enumi}{3}
\tightlist
\item
  Let us prove c). If \(a, b \in H\) , we have \(\theta_{a}\theta_{b}/\theta_{ab} \in K\) . It follows that the vector subspace of \(\mathrm{K}(\mathrm{H}^{1/n})\) generated by the \(\theta_{a}\) is stable under multiplication and so coincides with \(\mathrm{K}(\mathrm{H}^{1/n})\) . It only remains to show that the \(\theta_{a}\) are linearly independent; to do this we may evidently assume that \(H/K^{*n}\) is finite; then \([\mathrm{K}(\mathrm{H}^{1/n}): \mathrm{K}] = (\mathrm{Gal}(\mathrm{K}(\mathrm{H}^{1/n})/\mathrm{K}) : \{1\}) = (\mathrm{H}: \mathrm{K}^{*n})\) by b) and Cor. 2 of V, p.~87; since the number of \(\theta_{a}\) is equal to \((\mathrm{H}: \mathrm{K}^{*n})\) and they generate the vector K-space \(\mathrm{K}(\mathrm{H}^{1/n})\) , they are linearly independent.
\end{enumerate}

Examples. --- 1) There exists a largest abelian extension of exponent dividing n of K, contained in \(\Omega\) ; it is obtained by adjoining to K the n-th roots of all its elements; its Galois group may be identified with \(\operatorname{Hom}(\mathbf{K}^{*}/\mathbf{K}^{*n}, \mu_{n}(\mathbf{K}))\) , hence also with \(\operatorname{Hom}(\mathbf{K}^{*}, \mu_{n}(\mathbf{K}))\) .

* 2) Let us take \(\mathbf{K} = \mathbf{Q}\) and \(n = 2\) . Then \(\mathbf{Q}^{*}/\mathbf{Q}^{*2}\) is a vector \(\mathbf{F}_{2}\) -space having as a basis the union of \(\{-1\}\) and the set of all prime numbers. The largest abelian extension of exponent 2 of \(\mathbf{Q}\) contained in \(\mathbf{C}\) is thus the subfield \(\mathbf{Q}(i, \sqrt{2}, \sqrt{3}, \sqrt{5}, \ldots)\) of \(\mathbf{C}\) . Its Galois group consists of all the automorphisms obtained by multiplying each of the elements \(i\) , \(\sqrt{2}\) , \(\sqrt{3}\) , \(\sqrt{5}\) etc. by \(\pm 1\) . *

\begin{enumerate}
\def\labelenumi{\arabic{enumi})}
\setcounter{enumi}{2}
\item
  Let \(L\) be a cyclic extension of \(K\) of degree \(n\) ; then the group \((L^n \cap K^*) / K^{*n}\) is cyclic of order \(n\) . If \(a \in K^*\) is such that the residue class of \(a \mod K^{*n}\) is a generator of this group, then \(L\) is \(K\) -isomorphic to \(K[X] / (X^n - a)\) and the group \(\operatorname{Gal}(L / K)\) consists of the \(n\) automorphisms mapping \(X\) to \(\zeta X\) , \(\zeta \in \mu_n(K)\) .
\item
  Conversely let \(a \in \mathbf{K}^*\) and let \(r\) be the least integer \(> 0\) such that \(a^r \in \mathbf{K}^{*n}\) ; then the subfield \(L\) of \(\Omega\) generated by the roots of the polynomial \(X^n - a\) is a cyclic extension of \(K\) of degree \(r\) . In particular, \(X^n - a\) is irreducible if and only if \(r = n\) .
\end{enumerate}

Remark. --- Let \(a \in K^{*}\) and let r be the least integer \textgreater0 such that \(a^{r} \in K^{n}\) . Let B be the set of roots in K of the polynomial \(X^{n/r} - a\) ; then we have

\[
\mathbf {X} ^ {n} - a = \prod_ {b \in B} \left(\mathbf {X} ^ {r} - b\right),\tag{14}
\]

by substitution of \(X^{r}\) for T in the relation \(T^{n/r}-a=\Pi(T-b)\) . By Example 4 each of the polynomials \(X^{r}-b\) is irreducible, so that (14) is the decomposition of \(X^{n}-a\) into irreducible polynomials in K{[}X{]}.

